\exercisesection{}

\begin{enumerate}
\def\labelenumi{\arabic{enumi})}
\item
  设 \(G\) 为紧群。证明每个连续表示 \(\varphi\) 都是从 \(G\) 到 \(\mathbf{R}_+^*\) 的表示，并满足 \(\varphi(G) = \{1\}\)。由此推出 \(G\) 上的相对不变正测度是不变的。
\item
  2. 设 G 为拓扑群，且 G 的换位子群在 G 中稠密。证明 G 到 Hausdorff 交换群的每个连续表示 \(\varphi\) 都满足 \(\varphi(G) = \{e\}\)。由此推出，若 G 局部紧，则 G 上的每个相对不变复测度都是不变的。特别地，G 是幺模的。
\item
  设 \(G\) 为局部紧群，\(\mu\) 为 \(G\) 上的左 Haar 测度，且 \(\nu = \Delta_G^{-1/2} \cdot \mu\)。证明
\end{enumerate}

\[
\boldsymbol {\gamma} (s) \nu = \Delta_ {\mathrm{G}} (s) ^ {1 / 2} \nu , \quad \boldsymbol {\delta} (s) \nu = \Delta_ {\mathrm{G}} (s) ^ {1 / 2} \nu , \quad \stackrel {\vee} {\nu} = \nu .
\]

\begin{enumerate}
\def\labelenumi{\arabic{enumi})}
\setcounter{enumi}{3}
\item
  设 \(G, G'\) 为两个局部紧群，\(V\)（分别为 \(V'\)）为 \(G\)（分别为 \(G'\)）的单位元的一个开邻域，\(\varphi\) 为 \(G'\) 到 \(G\) 的局部同构，定义在 \(V'\) 上，且 \(\varphi(V') = V\)。证明 \(\Delta_G \circ \varphi\) 是 \(\Delta_{G'}\) 在 \(V'\) 上的限制。
\item
  对每个 \(a\)（属于乘法群 \(\mathbf{Q}^*\)），令 \(\varphi(a)\) 为加法群 \(\mathbf{R}\) 的自同构，定义为 \(\varphi(a)x = ax\)。赋予 \(\mathbf{Q}^*\) 离散拓扑。令 \(\mathbf{G}\) 为 \(\mathbf{Q}^*\) 与 \(\mathbf{R}\) 的由 \(\varphi\) 定义的拓扑半直积（GT, III, §2, No.~10）。证明 \(\mathbf{G}\) 局部同构于 \(\mathbf{R}\)，但不是幺模的。
\item
  设 \(\mathbf{G}\) 为含有一个开紧子群 \(\mathbf{H}\) 的局部紧群。证明对每个自同构 \(\varphi\)（它作用于 \(\mathbf{G}\)），有 mod \(\varphi \in \mathbf{Q}_+^*\)。（注意，\(\varphi(\mathbf{H}) \cap \mathbf{H}\) 在 \(\mathbf{H}\) 和 \(\varphi(\mathbf{H})\) 中的指数都是有限的。）证明等式 \(\Delta_{\mathbf{G}}(s) = 1\) 定义了 \(\mathbf{G}\) 的一个包含 \(\mathbf{H}\) 的开子群。
\item
  7. 设 G 为局部紧群，\(\beta\) 为 G 上的左 Haar 测度，A 为 G 的一个子集，B 为 G 的一个相对紧的 \(\beta\)-可积子集，且 \(\beta(\mathrm{B}) > 0\)。证明，若 \(\mathbf{A} \cdot \mathbf{B}\) 的紧子集对于 \(\beta\) 具有有界测度，则 \(\mathbf{A}\) 是相对紧的。（仿照命题 2 的论证。）
\item
  8. 设 G 为局部紧群，A 为 G 的稠密子集，\(\alpha\) 为 G 上的左 Haar 测度，H 为 G 的一个 \(\alpha\)-可测子集，具有如下性质：对每个 \(s \in A\)，\(sH \cap (\mathbb{C}H)\) 和 \(H \cap (\mathbb{C}sH)\) 都是局部 \(\alpha\)-零测的。证明 H 或者是局部 \(\alpha\)-零测的，或者其补集是局部 \(\alpha\)-零测的。（证明 \(\varphi_{H} \cdot \alpha\) 是左不变的。）
\item
  采用第 2 号定理 1 的证明中的记号。令 \(\beta\) 为 G 上满足 \(\beta(f_0) = 1\) 的唯一左 Haar 测度。证明，对每个 \(f \in \mathcal{K}(\mathbf{G})\)，\(\mathrm{I}_g(f)\) 趋于 \(\beta(f)\)，关于 \(\mathfrak{B}\)。（令 \(a\) 为映射 \(g \mapsto \mathrm{I}_g(f)\) 关于 \(\mathfrak{B}\) 的一个聚点。存在一个超滤 \(\mathfrak{U}\)，比 \(\mathfrak{B}\) 更细，使得 \(\mathrm{I}_g(f)\) 趋于 \(a\)，关于 \(\mathfrak{U}\)。另一方面，\(\mathrm{I}_g(f)\) 趋于 \(\beta(f)\)，关于 \(\mathfrak{U}\)。）
\item
  10. 设 G 为局部紧群，\(\mu\) 为 G 上的左 Haar 测度。证明每个 \(\mu\)-可积集 A 都包含于紧集的一个可数并中。（化归到 A 为开集的情形，并注意 G 中存在一个开子群，它是紧子集的可数并。）
\end{enumerate}

¶ 11) 设 G 为在无穷远处可数的非离散局部紧群，β 为 G 上的左 Haar 测度。证明 G 的单位元 e 的每个紧邻域 V 都包含 G 的一个正规子群 H，该子群是 β-零测的，且 G/H 是一个 Polish 局部紧群。（可如下进行：令 L 为由 V 生成的 G 的开子群；令 \(b_{1}, b_{2}, \ldots\) 为关于 L 的左陪集代表元。构造单位元 e 的对称邻域的递减序列 \((\mathrm{V}_{n})\)，使其满足：1) \(\mathrm{V}_{n}^{2} \subset \mathrm{V}_{n-1}\)；2) \(x\mathrm{V}_{n}x^{-1} \subset \mathrm{V}_{n-1}\) 对每个 \(x \in \mathrm{V}\) 成立；3) \(b_{i}\mathrm{V}_{n}b_{i}^{-1} \subset \mathrm{V}_{n-1}\) 对 \(1 \leqslant i \leqslant n\) 成立；4) \(\beta(\mathrm{V}_{n}) \leqslant 1/n\)。置 \(\mathrm{H} = \mathrm{V}_{1} \cap \mathrm{V}_{2} \cap \cdots\)。

令 \((f_{n})\) 为关于 G 的左一致结构一致连续的数值函数序列。证明可以构造 H，使得除此之外，每个 \(f_{n}\) 在 H 的陪集上为常值。（在上述构造中，取 \(V_{n}\) 使得 \(|f_{i}(x) - f_{i}(y)| \leqslant 1/n\) 当 \(x^{-1}y \in V_{n}\) 且 \(1 \leqslant i \leqslant n\) 时。）

令 \((g_n)\) 为 \(\beta\)-可积的数值函数序列，定义在 \(G\) 上。证明可以构造 \(H\)，使得除此之外，每个 \(g_n\) 几乎处处等于一个在 \(H\) 的陪集上为常值的函数。

\begin{enumerate}
\def\labelenumi{\arabic{enumi})}
\setcounter{enumi}{11}
\tightlist
\item
  令 \(\mathbf{K}\) 为非离散局部紧域，\(\mathbf{E}\) 为 \(\mathbf{K}\) 上的左拓扑向量空间。
\end{enumerate}

\begin{enumerate}
\def\labelenumi{\alph{enumi})}
\item
  a) 若 E 是 1 维的，则 E 同构于 \(\mathbf{K}_s\)（仿照 TVS, I, §2, No.~2 的命题 2 的证明，将绝对值替换为函数 mod \(_K\)）。
\item
  b) 若 E 是 n 维的，则 E 同构于 \(K_{s}^{n}\)（仿照 TVS, I, §2, No.~3 的定理 2 的证明）。
\item
  c) 假设 E 局部紧。令 F 为 E 的一个有限维 n 的线性子空间。~对每个 \(a \in K\)，令 \(\text{mod}_{\mathbf{E}}(a)\) 为 E 中映射 \(x \mapsto ax\) 的模。由于根据 b)，F 在 E 中是闭的，可以构造 \(\text{mod}_{\mathbf{E}/\mathbf{F}}(a)\)。证明 \(\text{mod}_{\mathbf{E}}(a) = \text{mod}_{\mathbf{K}}(a)^n \text{mod}_{\mathbf{E}/\mathbf{F}}(a)\)。证明，若 \(\text{mod}_{\mathbf{K}}(a) < 1\) 且 \(F \neq E\)，则 \(\text{mod}_{\mathbf{E}/\mathbf{F}}(a) < 1\)。（有 \(a^n \to 0\)，当 \(n \to +\infty\) 时；由此仿照 No.~10 的命题 12 的论证，推出 \(\text{mod}_{\mathbf{E}/\mathbf{F}}(a^n) \to 0\)。）由此推出 \(n \leqslant \log \text{mod}_{\mathbf{E}}(1/a)/\log \text{mod}_{\mathbf{K}}(1/a)\)，从而 E 是有限维的。
\end{enumerate}

（我们将在 CA, VI 中看到，局部紧域的拓扑可以由一个绝对值定义。于是 c) 的最终结果将成为 TVS, I, §2, No.~4 的定理 3 的特例。）

¶ 13) 设 G 为局部紧群，连续地左作用于局部紧 Polish 空间 T；β 为 G 上的左 Haar 测度，ν 为在 G 下拟不变的 T 上的正测度。

\begin{enumerate}
\def\labelenumi{\alph{enumi})}
\tightlist
\item
  存在函数 \((s,x)\mapsto \chi (s,x) > 0\)，定义在 \(\mathbf{G}\times \mathbf{T}\) 上，它是局部 \((\beta \otimes \nu)\)-可积的，并且对每个 \(s\in \mathbf{G}\)，函数 \(x\mapsto \chi (s^{-1},x)\) 局部 \(\nu\)-可积且满足 \(\gamma(s)\nu = \chi(s^{-1},\cdot)\cdot\nu\)。（证明 \((s,x)\mapsto(s,sx)\) 将 \(\beta\otimes\nu\) 变为等价测度，可利用 Ch. V, §5, No.~5 的定理 2 的准则 3)；令 \(\chi(s^{-1},x)d\beta(s)d\nu(x)\) 为此等价测度。对于 \(\varphi\in\mathcal{H}(\mathrm{G})\) 和 \(\psi\in\mathcal{H}(\mathrm{T})\)，
\end{enumerate}

\[
\int \varphi (s) d \beta (s) \int \psi (s x) d \nu (x) = \int \varphi (s) d \beta (s) \int \psi (x) \chi (s ^ {- 1}, x) d \nu (x),
\]

从而 \(\int \psi(sx) d\nu(x) = \int \psi(x)\chi(s^{-1}, x) d\nu(x)\) 除去一个局部 \(\beta\)-零测集 \(N(\psi)\) 外成立。接着应用 Ch. VI, §3, No.~1 的引理 1；然后在一个零测集上修改 \(\chi\)，该集为 \((\beta \otimes \nu)\)-零测集。

\begin{enumerate}
\def\labelenumi{\alph{enumi})}
\setcounter{enumi}{1}
\tightlist
\item
  证明对于任意的 \(s, t\)（属于 \(\mathbf{G}\)），有
\end{enumerate}

\[
\chi (s t, x) = \chi (s, t x) \chi (t, x)
\]

除去一个 \(\nu\)-零测的 \(x\) 值集合外成立（利用关系 \(\pmb{\gamma}(st)\nu = \pmb{\gamma}(s)(\pmb{\gamma}(t)\nu)\)。c) 证明 a) 中的函数 \(\chi\)（见 \(a\)）由一个局部 \((\beta \otimes \nu)\)-零测子集 \(\mathbf{G} \times \mathbf{T}\) 决定。

¶ 14) 设 T 为局部紧空间，U 为 T 上的一致结构，且存在一个围道 \(V_{0}\)，使得 \(V_{0}(t)\) 对所有 \(t \in T\) 都是紧的（GT, II, §4, Exer. 9）。另一方面，令 \(\Gamma\) 为 T 的同胚群，它关于一致结构 U 一致等度连续，并且存在 \(a \in T\)，其在 \(\Gamma\) 下的轨道稠密。证明 T 上存在一个正测度 \(\neq 0\)，它在 \(\Gamma\) 下不变，并且除一个常数因子外唯一。可按如下方式进行：

\(1^{\circ}\) 关于不变测度的存在性，采用第 2 号定理 1 的方法；先证明：若 K 是 T 的一个紧子集，U 是 \(a\) 的一个开邻域，则存在有限个元素 \(\sigma_{i} \in \Gamma\)，使得 \(\mathrm{K} \subset \bigcup \sigma_{i}(\mathrm{U})\)。

\(2^{\circ}\) 关于唯一性，首先注意，\(\mathcal{U}\) 的围道中在每个 \(\sigma \times \sigma\)（\(\sigma \in \Gamma\)）这个 T × T 的同胚下不变者，构成围道的一个基本系 \(\mathfrak{S}\)。对于 T 的每个相对紧集 A ⊂ T，以及每个具有非空内点的相对紧集 B ⊂ T，令 (A : B) 为用形如 \(\sigma\) B（\(\sigma \in \Gamma\)）的集合覆盖 A 所需的最少集合数；若 C ⊂ T 是第三个具有非空内点的相对紧集，则 (A : C) ≤ (A : B)(B : C)。令 K 为 T 的一个紧子集，令 L ⊃ K 为相对紧开集，令 V ∈ \(\mathfrak{S}\) 为包含于 V0 中的对称开围道，使得 V(K) ⊂ L，并令 W ∈ \(\mathfrak{S}\) 为包含于 V 中的对称闭围道。对每个满足 UW ⊂ V 且 WU ⊂ V 的对称围道 U ∈ \(\mathfrak{S}\)，以及每个 T 上的正测度 \(\nu\)，它在 \(\Gamma\) 下不变，有

\begin{enumerate}
\def\labelenumi{(\arabic{enumi})}
\tightlist
\item
\end{enumerate}

\[
\big (\mathrm{W} (a): \mathrm{U} (a) \big) \nu (\mathrm{K}) \leqslant \big (\mathrm{L}: \mathrm{U} (a) \big) \nu \big (\mathrm{V} (a) \big)\tag{2}
\]

\[
\bigl (\mathrm{K}: \mathrm{U} (a) \bigr) \nu \bigl (\mathrm{W} (a) \bigr) \leqslant \bigl (\mathrm{V} (a): \mathrm{U} (a) \bigr) \nu (\mathrm{L}).
\]

为此，证明：若 \((\sigma_i(\mathrm{U}(a)))\) 是由 \((\mathrm{L}:\mathrm{U}(a))\) 个集合构成的 L 的一个覆盖，则每个 \(x\in \mathbf{K}\) 至少属于 \((\mathrm{W}(a):\mathrm{U}(a))\) 个形如 \(\sigma_{i}(\mathrm{V}(a))\) 的集合，并且每个 \(y\in \mathrm{L}\) 至多属于 \((\mathrm{V}(a):\mathrm{U}(a))\) 个形如 \(\sigma_{i}(\mathrm{W}(a))\) 的集合（注意，若 \(y\in \sigma_i(\mathrm{W}(a))\)，则 \(\sigma_{i}(\mathrm{U}(a))\) 包含于 \(\mathrm{V}(x)\) 中）。

令 \(\mathfrak{F}\) 为一个比 \(\mathfrak{S}\) 的截面滤子更细的超滤；令 \(A_0\) 为 T 中一个非空的相对紧开集，并置 \(\lambda_U(A) = (A: U(a)) / (A_0: U(a))\)，其中 \(U \in \mathfrak{S}\) 为对称围道，A 为 T 的相对紧子集。令 \(\lambda(A) = \lim_{\mathfrak{F}} \lambda_U(A)\)；对 T 的每个紧子集 K，置 \(\lambda'(K) = \inf \lambda(B)\)，其中 B 遍历 K 的相对紧开邻域的集合。由 (1) 和 (2) 推出

\[
\begin{array}{r} \lambda^ {\prime} \big (\mathrm{W} (a) \big) \nu (\mathrm{K}) \leqslant \lambda (\mathrm{L}) \nu \big (\mathrm{W} (a) \big) \\ \lambda^ {\prime} (\mathrm{K}) \nu \big (\mathrm{W} (a) \big) \leqslant \lambda^ {\prime} \big (\mathrm{W} (a) \big) \nu (\mathrm{L}). \end{array}
\]

由此可知，若 \(K_{1}, K_{2}\) 是 T 的两个紧子集，且 \(L_{1} \supset K_{1}\)、\(L_{2} \supset K_{2}\) 是两个相对紧开集，则

\[
\lambda^ {\prime} (\mathrm{K} _ {1}) \nu (\mathrm{K} _ {2}) \leqslant \lambda (\mathrm{L} _ {2}) \nu (\mathrm{L} _ {1}),
\]

最后得到 \(\lambda^{\prime}(\mathrm{K}_{1})\nu (\mathrm{K}_{2}) = \lambda^{\prime}(\mathrm{K}_{2})\nu (\mathrm{K}_{1})\)

¶ 15) 设 G 为局部紧群，H 为 G 的闭子群，从而 G 连续地作用于 G/H。假设 H 满足如下条件：对 G 中 e 的每个邻域 V，存在 e 的一个邻域 U，使得 HU ⊂ VH（注意，如果 e 有一个在 G 的每个内自同构下不变的邻域基本系，则 G 的每个子群 H 都满足此条件）。证明 G/H 上存在一个在 G 下不变的非零正测度。（采用第 2 号定理 1 的证明中的相同方法，并注意：一方面，当 U 遍历 G 中 e 的邻域时，集合 HUH 在典范映射 π : G → G/H 下的像构成 π(H) 的邻域基本系；另一方面，对于 G 中 e 的每个邻域 V 和 G 的每个紧子集 K，存在 G 中 e 的一个邻域 U 和 G 的一个紧子集 L，使得每个与 KH 相交的形如 sHUH 的集合都可写成形如 s'HUH 的集合，其中 s' ∈ L。）

¶ 16) 设 G 为紧交换群，E 为 G 的稠密子集，且在 G 的运算下稳定。假设对 E 上的每个有界数值函数 f，都定义了一个数 M(f)，具有如下性质：对于 E 上任意两个有界函数 f、g，E 中任意元素 \(a_{i}, b_{k}\)（\(1 \leqslant i \leqslant m, 1 \leqslant k \leqslant n\)），以及任意实数 \(\alpha_{i}, \beta_{k}\)（\(1 \leqslant i \leqslant m, 1 \leqslant k \leqslant n\)），若

\[
\sum_ {i} \alpha_ {i} f (x a _ {i}) \leqslant \sum_ {k} \beta_ {k} g (x b _ {k})
\]

对所有 \(x \in \mathbf{E}\) 都有上述不等式，则有不等式

\[
\left(\sum_ {i} \alpha_ {i}\right) \mathrm{M} (f) \leqslant \left(\sum_ {k} \beta_ {k}\right) \mathrm{M} (g).
\]

另假设 \(\mathbf{M}(1) = 1\)（cf.~TVS, IV, Appendix）。

\begin{enumerate}
\def\labelenumi{\alph{enumi})}
\item
  令 \(\mu\) 为 G 上的归一化 Haar 测度。证明，对 G 上的每个连续数值函数 \(f\)，都有 \(\int f d\mu = \mathrm{M}(f|\mathrm{E})\)。（通过取适当的单位分解，先在 E 上考虑恒等于 \(\int f d\mu\) 的常值函数，再用形如 \(x \mapsto \sum_{i} \alpha_i f(xa_i)\) 的函数（其中 \(a_i \in \mathrm{E}\)）去逼近它。）
\item
  b) 由 a) 推出：若对 E 的每个子集 B 置 \(\nu(B) = M(\varphi_B)\)，则对 G 的每个开子集 U，有 \(\nu(U \cap E) \geqslant \mu(U)\)，对 G 的每个闭子集 F，有 \(\nu(F \cap E) \leqslant \mu(F)\)。对于 G 的每个 \(\mu\)-二次集 P（Ch. IV, §5, Exer. 17 d)），有 \(\nu(P \cap E) = \mu(P)\)。
\end{enumerate}

¶ 17) 设 T 为局部紧空间，S 为 T 的一个子集，并赋予其运算 \((x,y)\mapsto xy\)，使 S 成为一个幺半群 \(^{1}\)（不预先假定有单位元）；当 S 取 T 的诱导拓扑时，该运算在 S×S 上连续。再假设 S 的每个元素都是可消去的 \(^{2}\)（A, I, §2, No.~2）。令 \(\mu\) 为 T 上的有界正测度，集中在 S 上（因而 S 是 \(\mu\)-可测的），且总质量为 1。假设 \(\mu\) 在以下意义下左不变：对于 S 上定义的每个连续有界数值函数 f（根据 Ch. IV, §5, No.~5 的 Prop. 8 的 Cor.，因而 \(\mu\)-可测），都有 \(\int f(sx) d\mu(x) = \int f(x)d\mu(x)\)，对所有 \(s \in S\) 成立。这等价于说，对 S 的每个紧子集 K，都有 \(\mu(sK) = \mu(K)\)。

\begin{enumerate}
\def\labelenumi{\alph{enumi})}
\item
  a) 考虑 T×T 上的积测度 \(\mu\otimes\mu\)；证明存在 S 的紧子集 K，使得 K×K 在两个连续映射 \((x,y)\mapsto(x,xy)\) 和 \((x,y)\mapsto(xy,x)\) 下的像的测度任意接近 1（利用 Lebesgue--Fubini 定理）。由此推出这些像有一个公共点，并据此证明 S 有单位元（cf.~A, I, §2, Exer. 9）。
\item
  b) 证明 S 是紧群。（首先证明对每个 \(x \in S\)，内测度 \(\mu_{*}(xS)\) 等于 1，从而 xS 可测且测度集中在 xS 上；xS 随后成为一个可以应用 a) 的结果的幺半群，由此证明 S 是群。仿照命题 2 的论证证明 S 是紧的。最后利用 GT, III, §4.3 的 Exer. 21。）
\end{enumerate}

18) 设 X 为局部紧空间，G 为紧群，连续地作用于 X；E 为 X 中一个点的轨道，\(\mathcal {F}\) 为 X 上连续数值函数所成的向量空间，使得对每个函数 \(f \in \mathcal {F}\) 和每个 \(s \in G\)，都有 \(\gamma (s) f \in \mathcal {F}\)；再假设 \(\mathcal {F}\) 包含 X 上的常值函数。令 \(x _ {0} \in \mathrm{X}\) 为 G 下不变的点，并满足 \(| f (x _ {0}) | \leqslant \sup _ {y \in E} | f (y) |\) 对每个函数 \(f \in \mathcal {F}\) 成立。证明

从而有 \(f(x_0) = \int_{\mathbf{G}} f(s \cdot z) ds\)，对每个 \(z \in \mathbf{E}\) 和每个 \(f \in \mathcal{F}\) 成立。（证明存在一个正测度 \(\nu\)，定义在 \(\mathbf{E}\) 上，总质量为 1，使得 \(f(x_0) = \int_{\mathbf{E}} f(z) d\nu(z)\)，对每个 \(f \in \mathcal{F}\) 成立，然后应用 Lebesgue-Fubini 定理。）*当 \(X = R^n\)、\(G\) 为正交群且 \(x_0 = 0\) 时；应用 §2 的公式 (7)。*

\begin{enumerate}
\def\labelenumi{\arabic{enumi})}
\setcounter{enumi}{18}
\tightlist
\item
  设 X 为紧空间，A 为带单位元的 R 上赋范代数，G 为紧群；假设 G 连续地作用于 A 和 X，且对每个 \(s \in G\)，映射 \(a \mapsto s \cdot a\) 是代数 A 的自同构，并且满足 \(\|s \cdot a\| = \|a\|\) 对所有 \(a \in A\)。称从 X 到 A 的映射 f 在 G 下协变，如果
\end{enumerate}

\[
f (s \cdot x) = s \cdot f (x)
\]

对所有 \(s \in G\) 和 \(x \in X\) 都成立。令 \(B\) 为 \(A^X\) 的一个子环，由协变连续函数组成，并包含所有从 \(X\) 到 \(R\) 的协变连续映射（其中 \(R\) 被认同为 \(A\) 的子代数；注意，对这样的映射 \(g\)，有 \(g(s \cdot x) = g(x)\) 对所有 \(s \in G\) 和 \(x \in X\) 成立）。令 \(f\) 为从 \(X\) 到 \(A\) 的协变连续映射，并假设对每个 \(y \in X\)，存在一个映射 \(g_y \in B\)，使得 \(f(y) = g_y(y)\)。证明，对每个 \(\varepsilon > 0\)，存在一个映射 \(g \in B\)，使得 \(\|f(x) - g(x)\| \leqslant \varepsilon\) 对所有 \(x \in X\) 成立。（利用适当的连续单位分解（\(\varphi_i\)）（定义在 \(X\) 上）并引入函数 \(h_i\)，使 \(h_i(x) = \int_G \varphi_i(s \cdot x) ds\)。）

¶ 20) 设 G 为局部紧群，μ 为 G 上的左 Haar 测度，A 和 B 为 G 的两个子集。

\begin{enumerate}
\def\labelenumi{\alph{enumi})}
\tightlist
\item
  a) 假设下列两个条件之一成立：
\end{enumerate}

\(\alpha)\) A 是 \(\mu\)-可积的；

\(\beta)\ \mu^{*}(A) < +\infty\)，且 \(B\) 是 \(\mu\)-可测的。

证明，在这两种情形中的每一种情形下，函数 \(f(s) = \mu^{*}(sA \cap B)\) 关于 \(G\) 的右一致结构是一致连续的（关于 \(G\)）。(对于 \(M, N\) 为 \(G\) 的任意两个子集，置

\[
d (\mathrm{M}, \mathrm{N}) = \mu^ {*} \big ((\mathrm{M} \cap \mathbb {C N}) \cup (\mathrm{N} \cap \mathbb {C M}) \big).
\]

首先考虑 A 为紧集的情形；利用 Ch. IV, §4, No.~6, Th. 4，证明对每个 \(\varepsilon > 0\)，存在 G 中 \(e\) 的一个邻域 U，使得对每个 \(s \in G\) 和 \(t \in U\)，

\[
d (s \mathrm{A} \cap \mathrm{B}, s t \mathrm{A} \cap \mathrm{B}) \leqslant \varepsilon .
\]

然后应用 Ch. IV, §5 的 Exer. 13。若 B 是 \(\mu\)-可测的，\(\mu^{*}(\mathbf{A}) < +\infty\)，且 \((\mathbf{A}_{n})\) 是 G 的一个包含 A 的可积子集递减序列，满足 \(\inf (\mu (\mathbf{A}_n)) = \mu^{*}(\mathbf{A})\)，证明

\[
\mu^ {*} (s \mathrm{A} \cap \mathrm{B}) = \inf \left(\mu^ {*} (s \mathrm{A} _ {n} \cap \mathrm{B})\right)
\]

（Ch. V, 1st edn., §2, No.~2, Lemma 1）；\(^{4}\) 但注意另一方面，\(\mu(A_{n}-A_{n+1})\) 随 \(1/n\) 趋于 0。）

\begin{enumerate}
\def\labelenumi{\alph{enumi})}
\setcounter{enumi}{1}
\item
  若 \(\mathbf{A}^{-1}\) 是 \(\mu\)-可积的且 \(\mu^{*}(\mathbf{B}) < +\infty\)，则函数 \(f\) 关于 \(\mathbf{G}\) 的右、左一致结构都是一致连续的；此外，有 \(\int_{\mathbf{G}} f(s) d\mu(s) = \mu(\mathbf{A}^{-1})\mu^{*}(\mathbf{B})\)。（化归到 \(\mathbf{B}\) 可积的情形；注意此时 \(\mu(s\mathbf{A} \cap \mathbf{B}) = \mu(\mathbf{A} \cap s^{-1}\mathbf{B})\)，\(\varphi_{s\mathbf{A} \cap \mathbf{B}} = \varphi_{s\mathbf{A}}\varphi_{\mathbf{B}}\)，且 \(\varphi_{s\mathbf{A}}(t) = \varphi_{t\mathbf{A}^{-1}}(s)\)。）
\item
  由 \(a\) 推出，在上述两种情形下，只要 A 和 B 都不是零测的，AB 和 BA 的内部都非空。（Cf. Ch. VIII, §4, No.~6, Prop. 17。）
\item
  d) 在群 \(\mathbf{G} = \mathbf{SL}_{2}(\mathbf{R})\) 中，给出一个紧集 A 和一个 \(\mu\)-可测集 B 的例子，使得 \(f(s) = \mu(sA \cap B)\) 关于左一致结构不是一致连续的。（注意，存在 G 中趋于 e 的元素序列 \((t_{n})\) 和 G 中的元素序列 \((s_{n})\)，使得序列 \((s_{n}^{-1}t_{n}s_{n})\) 趋于无穷远点。）
\item
  e) 在局部紧群 G 中，给出两个不可测集 A、B 的例子，它们的外测度有限，且函数 \(s \mapsto \mu^{*}(sA \cap B)\) 不连续（cf. Ch. IV, §4, Exer. 8）。
\end{enumerate}

\begin{enumerate}
\def\labelenumi{\arabic{enumi})}
\setcounter{enumi}{20}
\tightlist
\item
  \begin{enumerate}
  \def\labelenumii{\alph{enumii})}
  \tightlist
  \item
    设 G 为局部紧群，\(\mu\) 为 G 上的左 Haar 测度。证明，若 S 是 G 的一个稳定子集且 \(\mu_{*}(S)>0\)，则 S 的内部非空（cf.~Exer. 20）。特别地，若 G 的一个子群 H 具有非零内测度，则 H 是 G 的开子群。
  \end{enumerate}
\end{enumerate}

\begin{enumerate}
\def\labelenumi{\alph{enumi})}
\setcounter{enumi}{1}
\item
  若 G 是紧的，则 G 的每个满足 \(\mu_{*}(S) > 0\) 的稳定子集 S 都是 G 的紧开子群（注意，S 的内部是稳定的，并利用 Exer. 17 b)）。
\item
若 G 是紧交换群，采用加法记号，并且 G 中存在一个无限阶元素 a。证明存在 G 的一个稳定子集 S，使 \(a \in S\)、\(-a \notin S\) 且 \(G = S \cup (-S)\)（利用 Zorn 引理）；由 b) 推出 S 不可测且 \(\mu_{*}(S) = 0\)。
\end{enumerate}

\begin{enumerate}
\def\labelenumi{\arabic{enumi})}
\setcounter{enumi}{21}
\tightlist
\item
  设 G 为局部紧群，\(\mu\) 为 G 上的左 Haar 测度，A 为 G 的一个可积子集，且 \(\mu(A) > 0\)。证明集合 H(A) 由满足 \(s \in G\) 且 \(\mu(A) = \mu(A \cap sA)\) 的元素组成，是一个紧群。（利用 Exer. 20，注意 H(A) 是 G 的闭子群。为证明其紧性，取 A 的一个紧子集 B，使 \(\mu(B) > \mu(A)/2\)，并证明 \(H(A) \subset BB^{-1}\)。）
\end{enumerate}

¶ 23) 设 G 为局部紧交换群，采用加法记号；μ 为 G 上的 Haar 测度，A、B 为两个可积子集。

\begin{enumerate}
\def\labelenumi{\alph{enumi})}
\tightlist
\item
  a) 对每个 \(s \in \mathbf{G}\)，令
\end{enumerate}

\[
\mathrm{A} ^ {\prime} = \sigma_ {s} (\mathrm{A}, \mathrm{B}) = \mathrm{A} \cup (\mathrm{B} + s), \quad \mathrm{B} ^ {\prime} = \tau_ {s} (\mathrm{A}, \mathrm{B}) = (\mathrm{A} - s) \cap \mathrm{B}.
\]

证明 \(\mu (\mathrm{A}^{\prime}) + \mu (\mathrm{B}^{\prime}) = \mu (\mathrm{A}) + \mu (\mathrm{B})\) 且 \(\mathrm{A}' + \mathrm{B}'\subset \mathrm{A} + \mathrm{B}\)

\begin{enumerate}
\def\labelenumi{\alph{enumi})}
\setcounter{enumi}{1}
\tightlist
\item
  假设 0 属于 \(\mathbf{A} \cap \mathbf{B}\)。称 \((\mathbf{A}', \mathbf{B}')\)（\(\mathbf{G}\) 的可积子集）由 (A, B) 导出，如果存在元素的序列 \((s_k)_{1 \leqslant k \leqslant n}\)，其元素属于 \(\mathbf{G}\)，和两个子集序列 \((\mathbf{A}_k)_{0 \leqslant k \leqslant n}\)、\((\mathbf{B}_k)_{0 \leqslant k \leqslant n}\)，它们属于 \(\mathbf{G}\)，使得 \(\mathbf{A}_0 = \mathbf{A}\)、\(\mathbf{B}_0 = \mathbf{B}\)、\(\mathbf{A}_k = \sigma_{s_k}(\mathbf{A}_{k-1}, \mathbf{B}_{k-1})\)、\(\mathbf{B}_k = \tau_{s_k}(\mathbf{A}_{k-1}, \mathbf{B}_{k-1})\) 对 \(1 \leqslant k \leqslant n\) 成立，\(s_k \in \mathbf{A}_{k-1}\) 对 \(1 \leqslant k \leqslant n\) 成立，且 \(\mathbf{A}' = \mathbf{A}_n\)、\(\mathbf{B}' = \mathbf{B}_n\)。证明存在一列有序对 \((\mathrm{E}_n, \mathrm{F}_n)\)，使得 \(\mathbf{E}_0 = \mathbf{A}\)、\(\mathbf{F}_0 = \mathbf{B}\)、\((\mathrm{E}_{n+1}, \mathrm{F}_{n+1})\) 由 \((\mathrm{E}_n, \mathrm{F}_n)\) 导出，并且 \(\mu((\mathrm{E}_n - s) \cap \mathrm{F}_n) \geqslant \mu(\mathrm{F}_{n+1}) - 2^{-n}\) 对每个 \(n\) 和每个 \(s \in \mathrm{E}_n\) 成立。置 \(\mathbf{E}_{\infty} = \bigcup_{n} \mathbf{E}_n\)、\(\mathbf{F}_{\infty} = \bigcap_{n} \mathbf{F}_n\)。证明对每个 \(s \in \mathrm{E}_{\infty}\)，
\end{enumerate}

\[
\mu \big ((\mathrm{E} _ {\infty} - s) \cap \mathrm{F} _ {\infty} \big) = \mu (\mathrm{F} _ {\infty}).
\]

\begin{enumerate}
\def\labelenumi{\alph{enumi})}
\setcounter{enumi}{2}
\tightlist
\item
  假设 \(\mu (\mathrm{F}_{\infty}) > 0\)。证明函数
\end{enumerate}

\[
f (s) = \mu \big ((\mathrm{E} _ {\infty} - s) \cap \mathrm{F} _ {\infty} \big)
\]

它只能取 0 和 \(\mu(\mathrm{F}_{\infty})\) 这两个值；并且集合 C 由满足 \(s \in G\) 且 \(f(s) = \mu(\mathrm{F}_{\infty})\) 的点组成，既开又闭，满足 \(\mu(C) = \mu(\mathrm{E}_{\infty})\)，且是 \(\mathrm{E}_{\infty}\) 的闭包（利用 Exer. 20 a) 和 b)）。另一方面，令 D 为满足 \(s \in \mathrm{F}_{\infty}\) 且 \(\mathrm{F}_{\infty}\) 与 \(s\) 的每个邻域的交的测度为 \(>0\) 的点所成的集合。证明 \(\mu(D) = \mu(\mathrm{F}_{\infty})\)，并且

\[
\mathbf {E} _ {\infty} + \mathbf {D} \subset \mathbf {C};
\]

由此推出 D 包含于 Exer. 22 所定义的子群 H(C)，且 H(C) 是 G 的紧开子群。最后证明 C + H(C) = C，\(\mu(C) \geqslant \mu(A) + \mu(B) - \mu(H(C))\)，以及 C ⊂ A + B（对每个 c ∈ C，考虑 E∞ ∩ (c - F∞) 的测度）。

\begin{enumerate}
\def\labelenumi{\alph{enumi})}
\setcounter{enumi}{3}
\tightlist
\item
  由 c) 推出，对于 G 的两个可积子集 A、B：或者 \(\mu_{*}(A+B)\geqslant\mu(A)+\mu(B)\)；或者存在 G 的一个紧开子群 H，使得 \(A+B\) 包含 H 的一个陪集，此时 \(\mu_{*}(A+B)\geqslant\mu(A)+\mu(B)-\mu(H)\)。考虑 G 连通的情形。
\end{enumerate}

¶ 24 a) 在 R 中，令 A（分别为 B）为具有如下性质的数 x 的集合：它们的正规二进展开 \(x = x_0 + \sum_{i=1}^{\infty} x_i 2^{-i}\)（\(x_0\) 为整数，\(x_i = 0\) 或 \(x_i = 1\)，当 \(i \geqslant 1\) 时）满足 \(x_{i}=0\) 对 i 为偶数且 \textgreater0（分别地，i 为奇数时满足此条件）。证明，对于 Lebesgue 测度，A 和 B 的测度为零，但 \(A+B=R\)。

\begin{enumerate}
\def\labelenumi{\alph{enumi})}
\setcounter{enumi}{1}
\item
  由 a) 推出，R 中存在一个包含于 A ∪ B 的基 H（在 Q 上），因而其测度为零。数字 rh（其中 r ∈ Q、h ∈ H）组成的集合 P₁ 也具有测度零。
\item
  令 \(P_{n}\) 为这样的实数集合：相对于基 H，其坐标中至多有 n 个非零坐标。证明，若 \(P_{n}\) 是零测的且 \(P_{n+1}\) 可测，则 \(P_{n+1}\) 是零测的。（令 \(h_{0} \in H\)；先证明，\(x \in P_{n+1}\) 中相对于 \(h_{0}\) 的坐标为 \(\neq 0\) 的那些点所成的集合 S 是零测的。利用 Exer. 20 证明，若 \(P_{n+1}\) 不是零测的，则存在两个不同的点 \(x', x''\)，它们属于 \(P_{n+1} \cap CS\)，使得 \((x' - x'') / h_{0}\) 是有理数，并由此导出矛盾。）
\item
  d) 由 a) 和 b) 推出，R 中存在两个零测集 C、D，使得 C + D 不可测。
\end{enumerate}

¶ 25) a) 令 f 为定义在 R 上的正数值函数，它关于 R 上的 Lebesgue 测度 μ 可积、有界且支集紧。令 γ = sup f(t)。

对每个 \(w \in \mathbf{R}\)，记 \(\mathrm{U}_f(w)\) 为 \(t \in \mathbf{R}\) 中满足 \(f(t) \geqslant w\) 的点所成的集合；置 \(\nu_f(w) = \mu^*(\mathrm{U}_f(w))\)。证明，对每个 \(\alpha > 1\)，

\[
\int_ {- \infty} ^ {+ \infty} f ^ {\alpha} (t) d t = \int_ {0} ^ {\gamma} \nu_ {f} (w) \alpha w ^ {\alpha - 1} d w.
\]

\begin{enumerate}
\def\labelenumi{\alph{enumi})}
\setcounter{enumi}{1}
\tightlist
\item
  令 g 为满足与 f 相同条件的另一个数值函数，并置 \(\delta = \sup_{t \in \mathbf{R}} g(t)\)。令 h 为定义在 \(R^{2}\) 上的函数：\(h(u,v) = f(u) + g(v)\) 当 \(f(u)g(v) \neq 0\) 时，否则 \(h(u,v) = 0\)；最后置 \(k(t) = \sup_{u+v=t} h(u,v)\)，于是 k 是正的、可积的、有界的且支集紧。证明，对每个 \(\alpha > 1\)，
\end{enumerate}

\[
\int_ {- \infty} ^ {+ \infty} k ^ {\alpha} (t) d t \geqslant (\gamma + \delta) ^ {\alpha} \left(\frac {1}{\gamma^ {\alpha}} \int_ {- \infty} ^ {+ \infty} f ^ {\alpha} (t) d t + \frac {1}{\delta^ {\alpha}} \int_ {- \infty} ^ {+ \infty} g ^ {\alpha} (t) d t\right).
\]

（注意，当 \(0 < w < 1\) 时，有 \(\mathrm{U}_k(\gamma w + \delta w) \supset \mathrm{U}_f(\gamma w) + \mathrm{U}_g(\delta w)\)，并利用 \(a\) 和 Exer. 23 d)。）

\begin{enumerate}
\def\labelenumi{\alph{enumi})}
\setcounter{enumi}{2}
\tightlist
\item
  c) 令 \(\mu_{n}\) 为 \(\mathbf{R}^{n}\) 上的 Lebesgue 测度，A、B 为 \(\mathbf{R}^{n}\) 的两个可积子集。证明（Brunn-Minkowski 不等式）
\end{enumerate}

\[
\left(\left(\mu_ {n}\right) _ {*} (\mathrm{A} + \mathrm{B})\right) ^ {1 / n} \geqslant \left(\mu_ {n} (\mathrm{A})\right) ^ {1 / n} + \left(\mu_ {n} (\mathrm{B})\right) ^ {1 / n}.
\]

（化归到 A 和 B 都是紧集的情形。然后利用 Exer. 23 d)、Lebesgue--Fubini 定理、b) 中证明的不等式以及 Hölder 不等式，对 n 作归纳论证。）

¶ 26) 设 G 为局部紧群，μ 为 G 上的左 Haar 测度，k 为满足 ≥1 的整数，A 为可积集。证明，对每个 ε \textgreater{} 0，存在 G 中 e 的一个邻域 U，具有如下性质：对于 U 中含有 k 个元素的每个有限子集 S，满足 sS ⊂ A 的 s ∈ G 所成的集合的测度至少为 (1 - ε)μ(A)。（化归到 A 为紧集的情形，并取 U 使得

\[
\mu (\mathrm{AU}) \leqslant \left(1 - \frac {\varepsilon}{k - 1}\right) \mu (\mathrm{A}).
\]

对 G 的每个有限子集 H，置 \(p(\mathrm{H}) = \operatorname{Card}(\mathrm{H})\)，并置 \(q(\mathrm{H}) = 1\)（当 \(\mathrm{H} \neq \emptyset\) 时），以及 \(q(\mathrm{H}) = 0\)（当 \(\mathrm{H} = \emptyset\) 时）；通过计算积分

\[
\int_ {\mathbf {G}} p (\mathrm{A} \cap s \mathrm{S}) d s \quad \text { and } \quad \int_ {\mathbf {G}} q (\mathrm{A} \cap s \mathrm{S}) d s  ,
\]

证明：若 \(h \leqslant k\)，令 \(\mathbf{M}_h\) 为 \(s \in G\) 中满足 \(\operatorname{Card}(\mathbf{A} \cap s\mathbf{S}) = h\) 的元素所成的集合，则

\[
\sum_ {h = 1} ^ {k} h \mu (\mathrm{M} _ {h}) = k \cdot \mu (\mathrm{A}) \quad \text { and } \quad \sum_ {h = 1} ^ {k} \mu (\mathrm{M} _ {h}) \leqslant \mu (\mathrm{AU}).
\]

由此推出 \(\sum_{h=1}^{k-1} h\mu(M_h) \leqslant k\varepsilon \cdot \mu(A)\)。

\begin{enumerate}
\def\labelenumi{\arabic{enumi})}
\setcounter{enumi}{26}
\tightlist
\item
  设 \(\mathbf{G}\) 为作用于集合 \(\mathbf{X}\) 上的群。称 \(\mathbf{P}\)（分别为 \(\mathbf{C}\)）为 \(\mathbf{X}\) 的子集，并将其称为 G-填充集（分别为 G-覆盖集），如果对每个 \(s \neq e\)（\(\mathbf{G}\) 中的），都有 \(s\mathbf{P} \cap \mathbf{P} = \emptyset\)（分别地，如果 \(\mathbf{X} = \bigcup_{s \in \mathbf{G}} s\mathbf{C}\)）。称同时为 G-填充集和 G-覆盖集的子集 \(\mathbf{P}\) 为 G-铺砌集。
\end{enumerate}

\begin{enumerate}
\def\labelenumi{\alph{enumi})}
\item
  假设 X 局部紧，G 可数，连续地作用于 X，并且 X 上存在一个在 G 下不变的非零正测度 \(\mu\)。令 P、C 为 \(\mu\)-可积的 G-填充集和 G-覆盖集。证明 \(\mu(C) \geqslant \mu(P)\)。（注意，\(\mu(C) \geqslant \sum_{s \in G} \mu(C \cap sP)\)。）
\item
  假设此外 X 上存在一个与 X 的拓扑相容的一致结构，并且有一个由 G 不变的开围道组成的基本系 S。另一方面，令 \(\Delta(G)\) 为 X 的所有可积 G-覆盖集 C 的 \(\mu(C)\) 的下确界。令 V 为属于 S 的一个围道，令 \(a \in X\) 满足 \(\mu(V(a)) > \Delta(G)\)；证明存在 \(s \in G\)，使得 \(s \neq e\) 且 \((a, sa) \in \overline{V} \circ V\)。
\item
  假设 X 是局部紧群，\(\mu\) 是左 Haar 测度，G 是 X 的一个可数子群，并通过左平移作用。沿用对 \(\Delta(G)\) 的定义（见 \(b\)），证明，若 A 是 X 的可积子集且 \(\mu(A) > \Delta(G)\)，则存在 \(s \in G \cap AA^{-1}\)，使 \(s \neq e\)。特殊情形：若 \(X = R^n\)，且 G 是秩为 \(n\) 的离散子群，则 \(\mathbf{R}^n: \Delta(G)\) 等于 G 关于 Z 的一个基相对于 \(\mathbf{R}^n\) 的典范基的行列式绝对值；若 A 是内部非空的对称闭凸集 \(^5\)（在 \(\mathbf{R}^n\) 中），且 \(\mu(A) \geqslant 2^n \Delta(G)\)，证明 \(A \cap G\) 中存在一个异于 0 的点（Minkowski 定理）。
\item
在 a) 的假设下，令 \(f\) 为一个 \(\geqslant 0\) 且 \(\mu\)-可积的定义在 \(X\) 上的函数。证明存在 \(a, b\) 属于 \(X\) 的两个点，使得 \(\mu(C) \sum_{s \in G} f(sa) \geqslant \int_X f(x) d\mu(x)\) 且 \(\mu (\mathrm{P})\sum_{s\in \mathbb{G}}f(sb)\leqslant \int_{\mathbb{X}}f(x)d\mu (x).\)（注意，若 \(g\) 是一个可积函数 \(\geqslant 0\)，E 是
\end{enumerate}

\(\mu\)-可积集 \(\mathbf{X}\) 中，则存在 \(c \in \mathbf{E}\)，使得 \(\int_{\mathbf{E}} g(x) d\mu(x) \leqslant g(c)\mu(\mathbf{E})\)，并存在 \(c' \in \mathbf{E}\)，使得 \(\int_{\mathbf{E}} g(x) d\mu(x) \geqslant g(c')\mu(\mathbf{E})\)。

¶ 28) 在 Exer. 27 a) 的假设下，再假设存在一个 μ-可积的 G-铺砌集 F。

\begin{enumerate}
\def\labelenumi{\alph{enumi})}
\tightlist
\item
对于定义在 X 上的每个 \(\mu\)-可积数值函数 \(f \geqslant 0\)，置 \(\widetilde{f}(x) = \sum_{s \in G} f(sx)\)，从而 \(\int_{F} \widetilde{f}(x) d\mu(x) = \int_{X} f(x) d\mu(x)\)（§2, No.~10, Prop. 15）。令 \((f_i)_{1 \leqslant i \leqslant n}\) 为数值函数 \(\geqslant 0\)、\(\mu\)-可积于 X 的一个族；当 \(i \neq j\) 时，置
\end{enumerate}

\[
m _ {i j} = \int_ {\mathbf {F}} \widetilde {f} _ {i} (x) \widetilde {f} _ {j} (x) d \mu (x),
\]

并令 \(c_{i} = \sup_{x\in X}\widetilde{f}_{i}(x)\)。证明至少存在一对不同的指标 \((i,j)\)，使得

\[
m _ {i j} \geqslant \frac {1}{n (n - 1) \mu (\mathrm{F})} \left(\left(\sum_ {i = 1} ^ {n} \int_ {\mathrm{X}} f _ {i} (x) d \mu (x)\right) ^ {2} - \mu (\mathrm{F}) \sum_ {i = 1} ^ {n} c _ {i} \int_ {\mathrm{X}} f _ {i} (x) d x\right).
\]

（利用 Cauchy-Schwarz 不等式从下方估计和 \(\sum_{i\neq j}m_{ij}\)。）由此推出：若 \((\mathrm{A}_i)_{1\leqslant i\leqslant n}\) 是 \(\mu\)-可积的 G-填充集的有限族，且

\[
\sum_ {i = 1} ^ {n} \mu (\mathrm{A} _ {i}) > \mu (\mathrm{F}),
\]

则存在一对不同的指标 \((i,j)\) 以及 \(s\in G\)，使得 \(\mu (\mathrm{A}_is\cap \mathrm{A}_j) > 0\)。b) 若 \(\mathbf{G}_0\) 是 \(G\) 的有限指数子群，且 \((G:\mathbf{G}_0) = h\)，且 \((s_1,\ldots ,s_h)\) 是 \(\mathbf{G}_0\) 在 \(G\) 中的右陪集代表元系，则 \(\mathrm{F_0} = \bigcup_{1\leqslant i\leqslant h}s_i\mathrm{F}\) 是一个 \(G_{0}\)-铺砌集。

\begin{enumerate}
\def\labelenumi{\alph{enumi})}
\setcounter{enumi}{2}
\tightlist
\item
  在 \(\mathbf{X} = \mathbf{R}^n\) 中，令 \(\mathbf{A}\) 为内部非空的对称闭凸集；令 \(u_i : (x_j) \mapsto \sum_{j=1}^{n} c_{ij} x_j\) 为 \(m\) 个定义在 \(\mathbf{X}\) 上的线性型，其整数系数为 \(c_{ij} (m < n)\)。
\end{enumerate}

证明，对每个整数 \(p \geqslant 1\) 以及每个数 \(r \geqslant 0\)，若满足 \(\mu(A)r^n \geqslant 2^n p^m\)，则存在一个异于 0 的点 \(x \in rA \cap Z^n\)，并且 \(u_i(x) \equiv 0 \pmod{p}\) 对 \(1 \leqslant i \leqslant m\) 成立（将 Minkowski 定理（Exer. 27 c)）应用于子群 \(G_0\)，它是 \(Z^n\) 的子群，由 \(z \in Z^n\) 中满足 \(u_i(z) \equiv 0 \pmod{p}\) 对 \(1 \leqslant i \leqslant m\) 成立的元素构成，并利用 \(b\)）。特殊情形：当 \(n = 2\)、\(m = 1\) 且 A 由 \(|x_1| \leqslant 1\)、\(|x_2| \leqslant 1\) 定义时（Thue 定理）。

¶ 29) a) 令 p 为素数；存在两个整数 a、b，使 \(a^{2}+b^{2}+1 \equiv 0 \pmod{p}\)（Alg., Ch. V, 1st edn., §11, No.~5, Cor. of Th. 3）。证明存在不全为零的整数 \(x_{1}, x_{2}, x_{3}, x_{4}\)，使 \(ax_{1} + bx_{2} \equiv x_{3} \pmod{p}\)、\(bx_{1} - ax_{2} \equiv x_{4} \pmod{p}\)，并且

\[
y = x _ {1} ^ {2} + x _ {2} ^ {2} + x _ {3} ^ {2} + x _ {4} ^ {2} \leqslant \sqrt {2} p
\]

（采用 Exer. 28 c) 中的相同方法。）证明 \(y\) 能被 \(p\) 整除，并由此推出 \(y = p\)。

\begin{enumerate}
\def\labelenumi{\alph{enumi})}
\setcounter{enumi}{1}
\tightlist
\item
  b) 由 a) 推出，每个整数 \(n \geqslant 0\) 都是四个平方数之和（Lagrange 定理；利用 Alg., Ch. IV, 1st edn., §2, Exer. 11）。
\end{enumerate}

\begin{enumerate}
\def\labelenumi{\arabic{enumi})}
\setcounter{enumi}{29}
\tightlist
\item
  设 G 为局部紧群，\(\mu\) 为 G 上的左 Haar 测度，\(\nu\) 为 G 上的有界测度。假设映射 \(s \mapsto \gamma(s)\nu\) 从 G 到 Banach 空间 \(\mathcal{M}^{1}(G)\) 连续。证明测度 \(\nu\) 以 \(\mu\) 为基。（令 A 为 G 的一个 \(\mu\)-零测紧子集。仿照命题 11 的论证，证明对稠密子集中的 x 有 \(\nu(xA)=0\)。由此推出 \(\nu(A)=0\)。）
\end{enumerate}

反过来，若 \(\nu\) 以 \(\mu\) 为基，则映射 \(s \mapsto \pmb{\gamma}(s)\nu\) 从 G 到 \(\mathcal{M}^1(\mathrm{G})\) 连续：cf.~Ch. VIII, §2, No.~5, Prop. 8。
% semantic-fragments: legacy-input-seam
\relax{}
